\chapter{Методы}

\section{Спектральный анализ}

Основным примером для экспериментальной части этого исследования является методология \ref{alg:architecture_spectral}, которая была реализована в ноутбуке Resnet18-spectral-imagenet.ipynb. Данный подход был воспроизведен для CNN архитектуры, и адаптирован для анализа архитектур ViT и Swin.

\begin{algorithm}[h]
\caption{Обобщенная процедура вычисления спектральной характеристики архитектуры}
\label{alg:architecture_spectral}
\vspace{1em}
\begin{algorithmic}[1]
\Require Выбранный слой нейронной сети, геометрические параметры $H, W$.
\Ensure Нормализованная радиальная кривая спектра мощности $\tilde E(\rho)$.
\vspace{1em}

\State Получить представление $\mathcal{F}$ из выбранного слоя/блока
\If{$\mathcal{F}$ является 4D тензором (CNN)}
    \State Оставить тензор без изменений: $(B, H, W, D)$
\Else{ \ $\mathcal{F}$ является 3D тензором вида $(B, N, D)$ (ViT, SWIN)}
    \State Выполнить изменение формы в $(B, H, W, D)$ с $H=W=\sqrt{N}$
\EndIf
\State Центрировать тензор по пространству
\State Выполнить 2D FFT по осям $H,W$ для каждого канала и сдвиг центра 
\State Вычислить спектр мощности: $S(b, u, v, c) \gets |FFT(\dots)|^2$
\State Усреднить спектр по батчу и каналам: $\bar{S}(u, v)$
\For{Для каждого нормализованного радиуса $\rho = \frac{\sqrt{u^2+v^2}}{R_{max}}$ от $0$ до $1$, где $R_{max} = min(H,W)/2$}
    \State Найти все точки $(u, v)$, удовлетворяющие условию $\lfloor \sqrt{u^2 + v^2} \rfloor = r$
    \State Вычислить среднее значение $\bar{S}(u, v)$ по этим точкам и записать в $E(\rho)$
\EndFor
\State Нормализовать кривую: $E(\rho) \gets E(\rho) / \sum_{\rho} E(\rho)$
\State \Return $E(\rho)$
\end{algorithmic}
\end{algorithm}

Первоочередным действием является процедура извлечения представлений. В данной работе «представление» — это набор значений в некотором слое модели, организованный в пространственную структуру. Спектральный анализ опирается на предположение, что эти активации можно интерпретировать как дискретную функцию на двумерной решётке.

Для каждой из исследуемых моделей были реализованы функции перехвата (hook) на выходах ключевых блоков: layer1...4 для ResNet и равномерно по глубине для трансформеров ViT и SWIN. На каждый такой блок подаётся один и тот же набор изображений из валидационной выборки ImageNet \cite{priyerana_imagenet10k} (после стандартной нормализации). 
Реализация функций перехвата, отличается в зависимости от используемой архитектуры, поскольку отличаются представления, с которыми оперируют модели. Так, например в ResNet каждый слой явно оперирует с 2D картой признаков - тензоры имеют размерность $(B, C, H, W)$, где $H$ и $W$— пространственные размеры карты признаков на данном слое, $C$ — число каналов, $B$ — размер батча. 

ViT работает иначе: изображение разбивается на непересекающиеся патчи размера $16\times16$. Каждый патч линеаризуется в вектор и пропускается через линейный слой, превращаясь в токен размерности $D$. В результате получается последовательность токенов, длина которой составляет $N = W_{patch}\times H_{patch}$, где $W_{patch}, H_{patch}$ - количество патчей по горизонтали и вертикали. Внутри самой модели перед первым блоком к токенам добавляются позиционные эмбеддинги - таким образом сохраняется понимание пространственной структуры, итого на вход поступает $N+1$ токенов. Сами же токены обрабатываются через механизм самовнимания (self-attention), который не учитывает пространственную близость. Таким образом на выходе любого блока мы получаем не 2D карту, а одномерную последовательность размерности $(B, N+1,D)$ - с фиксированным порядком токенов, но внутри самих представлений отсутствует понимание соседства по вертикали или горизонтали. 

Поэтому для проведения с ViT аналогичного как у CNN спектральным анализом, в функции перехвата преобразуем последовательность токенов в 2D карту признаков. Исключаем из получаемой последовательности первый CLS-токен - который представляет собой пространственные эмбендинги. Находим размеры решетки патчей $H_{patch} = W_{patch} = \sqrt{N}$, поскольку после стандартной нормализации все изображения квадратные размера $224\times 224$. Затем преобразуем тензор из формы $(B, N, D)$ в $(B, C, H_{patch}, W_{patch})$ построчной разверткой таким образом, что $i$-й токен соответствует позиции $(i\mod W_p ,\lfloor i / W_p \rfloor)$. Однако есть важный нюанс, что эта карта признаков не эквивалентна карте признаков CNN. Поэтому пространственное разрешение в ViT равно количеству патчей, а не пикселей поскольку каждый токен - неперекрывающийся блок пикселей представляющий агрегированный вектор, не представляющий внутренней структуры.

В архитектуре SWIN разбиение входящего изображения происходит на непересекающиеся патчи размера $4\times4$. Это похоже на CNN тем, что также происходит уменьшение пространственного разрешения по мере углубления, так как на каждом этапе (stage) происходит слияние патчей, которое увеличивает количество каналов и уменьшает пространственное разрешение вдвое. Также в SWIN отсутствует CLS-токен, и преобразование в 2D карту признаков осуществляется аналогично с ViT.  Также отмечу, что внутренний механизм архитектуры смещение окон (shifted window) - не влияет на выход, и не требует специальной обработки для спектрального анализа.


Перед применением каких-либо частотных преобразований необходима предобработка.
Очень важным действием является вычитание пространственного среднего из каждой карты признаков:

\begin{equation}
    \tilde{x}_{b,c}(h,w) = x_{b,c}(h,w) - \frac{1}{HW} \sum_{h,w} x_{b,c}(h,w)
\end{equation}

Эта операция необходима для устранения нулевой частоты (постоянную составляющую , DC-компоненты) из сигнала. DC-компонента представляет из себя среднее значение по пространству и по абсолютной величине часто на порядки превосходит энергию всех остальных частот, тем самым маскируя вклад более высоких частот. Кроме того, постоянная составляющая несёт в себе информацию, которая для целей нашего исследования является скорее помехой.Её удаление необходимо сделать чтобы анализ был чувствительным именно к вариациям признаков, а не к их абсолютному уровню, для корректного анализа спектра мощности.
В рамках данной работы операция выполняется независимо для каждого канала и каждого образца в батче: из каждого элемента двумерной матрицы вычитается среднее арифметическое по всем элементам этой же матрицы. Альтернативой могло бы быть высокочастотное фильтрование, но оно требует подбора граничных частот и вносит дополнительный произвол; вычитание же среднего — это простейший и в то же время строгий способ обнуления DC-компоненты без искажения остального спектра.

Для перехода из пространственной области карт признаков в частотную используем двумерное дискретное преобразование Фурье, реализуемое через быстрого преобразования Фурье (БПФ). Получаем комплекснозначный спектр $F_{b,c}(u,v)$. 

\begin{equation}
    F_{b,c}(u, v) = \sum_{h=0}^{H-1} \sum_{w=0}^{W-1} x_{b,c}(h, w) e^{-2\pi i \left( \frac{uh}{H} + \frac{vw}{W} \right)}
\end{equation}
где $x(h,w)$ — значение признака в точке с координатами $h$ (высота) и $w$ (ширина), а $F(u,v)$ — комплексный коэффициент на частоте с индексами $u$ и $v$.
В результате для каждой карты признаков мы получаем комплексный массив того же пространственного размера, где действительная и мнимая части кодируют амплитуду и фазу соответствующих частотных компонент. Фаза, несёт информацию о пространственном расположении структур, но наша задача — анализ распределения энергии в низких или высоких частотах - в ней она не требуется. Поэтому на следующем этапе мы переходим к спектру мощности, отбрасывая фазовые соотношения.

Вычисляется спектр мощности (энергии) как квадрат модуля:

\begin{equation}
    S_{b,c}(u, v) = |F_{b,c}(u, v)|^2
\end{equation}

Энергетический спектр показывает, какой вклад вносит каждая пространственная частота в общую энергию сигнала, и инвариантен к фазе, что удобно для последующего усреднения.

Чтобы получить единую репрезентативную характеристику слоя, для отражения поведения в целом, применяем усреднение по всем образцам в батче и по всем каналам.

\begin{equation}
    \bar{S}(u, v) = \frac{1}{B \cdot C} \sum_{b=1}^{B} \sum_{c=1}^{C} S_{b,c}(u, v)
\end{equation}

Здесь $B$ - размер батча, в нашем исследовании это - 32, достаточное для устойчивой оценки, $C$ - число каналов на слое. Такое двойное усреднение позволяет, во-первых, нивелировать шум, связанный с конкретными примерами, а во-вторых, получить агрегированное представление о том, как слой в целом преобразует частотные компоненты. Важно отметить, что усреднение спектров мощности — это не то же самое, что усреднение самих карт признаков с последующим вычислением спектра. Это важно, поскольку на одних изображениям могут преобладать разнородные структуры - горизонтальные, а на других вертикальные, усреднение карт признаков привел бы к взаимному уничтожению этих разнородных признаков. 
Подытожив, усреднение энергетического спектра по каналам и батчам даёт единую спектральную карту для данного слоя, агрегированную по всему разнообразию признаков, которые сеть извлекла из обработанных изображений. 

Полученный после всех предыдущих шагов двумерный спектр $S(u,v)$ содержит информацию о зависимости энергии от направления частотного вектора. Однако для целей нашего исследования — анализа того, как меняется распределение энергии по мере увеличения пространственной частоты, вне зависимости от того, горизонтальная это частота или вертикальная — угловая зависимость является избыточной. Чтобы перейти к компактному одномерному представлению, мы применяем радиальное усреднение. Поскольку интерес представляет распределение энергии по пространственным частотам независимо от направления, выполняется переход к полярным координатам. 

Вычисляется среднее значение $\bar S(u,v)$ по всем точкам, попадающим в кольцо с внутренним радиусом $r$ и внешним радиусом $r+1$ (в дискретном случае). Поскольку частоты дискретны, радиус $r$ пробегает целые значения от $0$ до $R_{max}$. 
Полученная функция $E(r)$ зависит только от радиуса — то есть от абсолютной величины пространственной частоты. Для удобства интерпретации и сравнения слоёв с разным разрешением мы затем нормализуем радиус, деля его на $R_{max}$, что даёт безразмерную величину $\rho$. То есть для каждого нормализованного радиуса:

\begin{equation}
    \rho = \frac{\sqrt{u^2 + v^2}}{R_{\max}}, где \ \ R_{\max} = \frac{\sqrt{H^2 + W^2}}{2}
\end{equation}

Получаем одномерную функцию $E(\rho)$ определенную на интервале [0,1], которая несёт абсолютные значения энергии, усреднённой по всем направлениям. Где $\rho=0$ соответствует постоянной составляющей (после вычитания среднего она обнулена, так что значение $E(0)$ близко к нулю), а $\rho=1$ — частоте Найквиста, максимальной частоте, которую может нести карта признаков данного разрешения.

Однако сравнение таких функций для разных слоёв непосредственно по шкале $E(\rho)$ часто оказывается затруднительным поскольку абсолютная энергия может кардинально различаться в зависимости от глубины. Например, в глубоких слоях ResNet количество каналов растёт (с 64 до 2048), а пространственное разрешение падает. Суммарная энергия карт признаков может изменяться на порядки из-за масштабирования, вносимого нормализацией, функциями активации и остаточными связями. 

Для корректного сравнения слоёв с разной абсолютной энергией производится нормализация:

\begin{equation}
    \tilde{E}(\rho) = \frac{E(\rho)}{\sum_{\rho} E(\rho)}
\end{equation}

Сумма в знаменателе — это полная энергия всех частотных компонент (за вычетом DC, который мы обнулили на этапе предобработки). 
Нормализация превращает спектр в плотность распределения вероятности энергии по частотам. Теперь $\tilde E(\rho)$ можно интерпретировать как долю энергии, приходящуюся на окрестность нормированной частоты $\rho$.
У данного подхода есть нюанс: при нормализации мы теряем информацию об общей энергии слоя. В работе рассматриваются нормализованные спектры для качественного анализа - смещение пика в сторону низких частот, островершинное или плоское распределение. Также для количественной оценки вводятся дополнительные метрики, которые вычисляются на основе $\tilde E(\rho)$, такие как спектральный центроид и доля высоких и низких частот.

Слои в ResNet и Swin имеют разное пространственное разрешение. Чтобы корректно сравнивать их спектры, применяется масштабирование частотной оси. При сравнении кривых спектральной энергии для моделей одного семейства за эталон принимается разрешение первого слоя, и для всех остальных слоёв частотная шкала пересчитывается таким образом, чтобы физические частоты (в циклах на изображение) были сопоставимы. Это позволяет строить графики, на которых видно, как меняется распределение энергии по одним и тем же пространственным частотам по мере углубления в сеть. 

Для сравнения того как ведут себя кривые для разных архитектур, использую единое нормализованное частотное масштабирование к layer1 ResNet:

\begin{equation}
    \rho_{mapped} = \rho \cdot \frac{res_{layer}}{res_{ref}}  
\end{equation}

где $\text{res}_\text{ref}$ — разрешение эталонного слоя ( layer1 ResNet, поскольку имеет самое большое разрешение = 40). После масштабирования кривые имеют разные диапазоны по оси $x$. Для единообразного вычисления метрик выполняется интерполяция каждой кривой на общую сетку частот $[0, 1]$ с фиксированным шагом (например, 0.001). Это позволяет получить сопоставимые значения центроида и долей энергии. Это необходимо, потому что исходные спектры имеют разное количество точек (40, 20, 10, 5) и их прямое сравнение невозможно. Радиальные спектры $E(\rho)$ являются гладкими функциями поэтому линейная интерполяция даёт малую ошибку и не вносит осцилляций.
 
\section{Метрики}

\textbf{Спектральный центроид}

Спектральный центроид представляет собой средневзвешенную частоту, где весом выступает нормализованная энергия. Для дискретной радиальной спектральной кривой $E(\rho_k)$ центроид вычисляется как:

\begin{equation}
    C = \frac{\sum \rho_k \cdot E(\rho_k)}{\sum E(\rho_k)}
\end{equation}

Центроид принимает значения на отрезке [0,1] и показывает, где сосредоточена основная энергия. Низкое значение центроида свидетельствует о том, что энергия представления сильнее сосредоточена в низкочастотной области - представление сглажено и содержит крупномасштабные структуры; чем выше значение, тем больше вклад средних и высоких пространственных частот - в представлении много мелких деталей, текстур и резкие границы.

Эта метрика хороша инвариантностью к масштабу, что позволяет сравнивать между слоями с разным пространственным разрешением и между разными архитектурами. Также она чувствительна к изменениям - даже небольшие смещения спектра заметно меняют центроид, также она удобна для вычисления другой метрики - скорости сжатия.

\textbf{Доля энергии в низких и высоких частотах}

Помимо среднего, можно численно измерить, как энергия распределена в хвостах распределения. Для этого выбирается пороговое значение $\tau$ (например, $\tau = 0.2$ для низких частот). Доля энергии в низких частотах:
\begin{equation}
    LowFrac = \sum_{k:\rho_k<0.2} E(\rho_k)
\end{equation}

Аналогично для высоких частот (например, $\rho_k > 0.5$):
\begin{equation}
    HighFrac = \sum_{k:\rho_k>0.5} E(\rho_k)
\end{equation}

Эти метрики характеризуют, какая часть информации сохраняется в определённом частотном диапазоне. Рост LowFrac с глубиной означает, что сеть фокусируется на глобальных, контекстных признаках. Падение HighFrac до нуля говорит о полном подавлении мелких деталей. В экспериментах исследования HighFrac обнуляется уже на втором слое (для ResNet) или стадии (для SWIN) – это количественное подтверждение гипотезы о ранней потере высокочастотной информации.

Использование одного порога не различает низкие и средние частоты. Два порога позволяют сегментировать спектр на три содержательные области: низкие, средние и высокие частоты. Это даёт более детальную картину того, какие именно частоты сохраняются или теряются. Пороги 0.2 и 0.5 выбраны случайно, но могут быть адаптированы под конкретную задачу 

\textbf{Скорость «сжатия» спектра по глубине}

Также в качестве метрики можно построить зависимость центроида (или доли высоких частот) от индекса слоя (или относительной глубины), назовем метрику Спектральное сжатие. Спектральное сжатие – это процесс последовательного подавления высоких частот по мере углубления сети. Для его количественной оценки можно использовать линейную или экспоненциальную модель. 

В случае линейной модели предполагается, что центроид убывает линейно с глубиной:

\begin{equation}
    C_{i}= a \cdot i + b
\end{equation}

где $i$ – индекс слоя (или относительная глубина $i/N$). Коэффициент $a$ (отрицательный) – линейная скорость сжатия показывает на сколько в среднем уменьшается центроид при переходе к следующему слою.

Многие процессы в нейросетях носят экспоненциальный характер, также центроид убывает быстрее на первых слоях и замедляется к концу, что ближе к экспоненте, чем к прямой. Поэтому динамику лучше опишет экспоненциальная модель:

\begin{equation}
    C_{i}=C_{0} \cdot e^{-\lambda \cdot i}
\end{equation}

где $\lambda>0$ – экспоненциальная скорость затухания. Эта величина безразмерна. За один слой центроид уменьшается в $e^{-\lambda} $ раз. 

Другие метрики на подобие энтропии спектра или пиковой частоты не использовались так как неинформативны для данного исследования. Энтропия спектра характеризует «размытость», но не указывает, в какую сторону смещён спектр, поэтому она более подходит для исследования робастности. Пиковая частота - игнорирует форму распределения, не дает информации о широте спектра и больше чувствительна шуму.

Все метрики вычисляются для каждого батча независимо, после чего усредняются. Для каждой метрики $M$ сохранены не только среднее $\overline{M}$, но и стандартное отклонение $\sigma_M$ по батчам. Для визуального удобства в таблицах будут представлены только средние значения. 